\section{Нижняя оценка для произвольных запросов}\label{sec:lower}

Оценим снизу число запросов, необходимое любому методу,
который решает все задачи рассматриваемого класса.

\subsection{Происхождение четырёх слагаемых}\label{sec:lower-origin}

Нижнюю оценку докажем отдельно для каждого из четырёх слагаемых. Обозначим
\begin{equation}\label{eq:lower-four-terms}
 \begin{gathered}
 B_n=n,\\
 B_S=\sqrt{\frac{nL_0R^2}{\varepsilon}},\\
 B_Q=n(L_1R)^{2/3},\\
 B_X=\sqrt{\frac{nL_1^2R^2\delta}{\varepsilon}+n(L_1R)^2}.
 \end{gathered}
\end{equation}
Поскольку $\sqrt{a^2+b^2}\sim a+b$ при $a,b\ge0$, имеем
$T_0\sim B_n+B_S+B_X+B_Q$.

\begin{center}
\small
\renewcommand{\arraystretch}{1.18}
\begin{tabularx}{\linewidth}{@{}lYY@{}}
\toprule
Член & Что затрудняет решение & Трудное семейство\\
\midrule
$B_n$ & Поиск информативной компоненты & Скрытая пара функций\\
$B_S$ & Последовательное раскрытие гладких цепочек & Независимые цепочки S\\
$B_X$ & Неоднородность компонент при обобщённой гладкости & Экспоненциальные цепочки X\\
$B_Q$ & Повторный поиск компоненты для следующего шага & Цепочка Q со скрытым размещением звеньев\\
\bottomrule
\end{tabularx}
\end{center}

Для каждого слагаемого построим своё семейство задач.
Поскольку метод должен решать задачи всех четырёх семейств,
его сложность не меньше максимума полученных оценок.
Ниже зависимости обсуждаются с точностью до числовых и логарифмических
множителей. Точные условия конструкций приведены в приложении;
если отдельная конструкция неприменима, соответствующий член
оценивается через остальные.

\paragraph{Поиск информативной компоненты: $B_n=n$.}
Среди $n$ компонент случайно выбираются две ненулевые, остальные
равны нулю. До первого
обращения к одной из них алгоритм получает только нулевые ответы.
Эти ответы не позволяют различить две задачи с точками минимума по разные
стороны от начальной точки. Чтобы ожидаемая ошибка была мала,
нужно с достаточной вероятностью найти скрытую пару функций, а для этого
требуется порядка $n$ обращений к оракулу
(лемма~\ref{lem:lower-pair}).

\paragraph{Независимые гладкие цепочки: $B_S$.}
В каждую компоненту помещается своя цепочка координат. Под цепочкой
понимается функция, у которой очередная координата влияет на ответ
оракула только после достаточного продвижения по предыдущим
координатам. В идеализированной модели один запрос раскрывает не более
одного нового звена одной цепочки. Как это свойство переносится на
произвольные запросы, объяснено ниже.

В семействе S число звеньев одной цепочки имеет порядок
$\sqrt{L_0R^2/(n\varepsilon)}$. Чтобы достичь точности $\varepsilon$, нужно
раскрыть достаточную долю от общего числа звеньев всех цепочек. Поэтому
число запросов имеет порядок
\[
 n\sqrt{\frac{L_0R^2}{n\varepsilon}}
 =\sqrt{\frac{nL_0R^2}{\varepsilon}}=B_S.
\]
Задачи этого семейства гладкие: $\|\nabla^2f_i\|\le L_0$
(лемма~\ref{lem:lower-S}). Если условия построения цепочки не выполнены, то
$B_S\lesssim B_n$ и достаточно оценки для $B_n$.

\paragraph{Экспоненциальные цепочки: $B_X$.}
При обобщённой гладкости кривизна компоненты может расти вместе
с её градиентом. Даже в точке минимума среднего, где
$\nabla f(x^\ast)=0$, отдельные градиенты $\nabla f_i(x^\ast)$
могут быть ненулевыми и компенсировать друг друга. Поэтому малость нормы градиента среднего не влечёт
малости констант гладкости компонент.
Параметр $\delta$ ограничивает расхождение между минимумом среднего
и средним значением инфимумов компонент.

В семействе X цепочки входят в показатели
экспонент у компонент, объединённых в пары. Конструкция сохраняет последовательное
раскрытие звеньев и задаёт неоднородность, равную $\delta$.
Оценка для этого семейства вместе с оценкой для скрытой пары функций даёт $B_X$
(леммы~\ref{lem:lower-X} и~\ref{lem:lower-pair}). Поскольку
$B_X=\sqrt n\,L_1R\sqrt{1+\delta/\varepsilon}$, имеем
\[
  B_X\sim\sqrt n\,L_1R\max\left\{1,\sqrt{\frac\delta\varepsilon}\right\}.
\]

\paragraph{Поиск следующего звена: $B_Q$.}
В семействе Q используется одна цепочка, но её звенья случайно
размещены по парам компонент. После раскрытия очередной координаты
недостаточно знать, в каком направлении продолжать: нужно ещё
найти пару, в которой находится следующее звено. Компонента, которая уже
была запрошена, может стать информативной лишь на следующем уровне
цепочки, поэтому однократное вычисление всех компонент не избавляет
от этих поисков.

При условиях леммы~\ref{lem:lower-Q} длина цепочки имеет порядок
$(L_1R)^{2/3}$, а поиск пары на очередном уровне требует порядка $n$
запросов в среднем. Отсюда с точностью до логарифмических множителей
\[
 \underbrace{(L_1R)^{2/3}}_{\text{число уровней}}
 \;\underbrace{n}_{\text{поиск на уровне}}
 =B_Q.
\]
Показатель $2/3$ объясняется выбором длины цепочки $T$.
В точке минимума её $T$ активных координат равны $R/\sqrt T$, так
что норма этой точки равна $R$. Оценка невязки через число раскрытых
звеньев~\eqref{eq:lower-Q-gap} показывает, что для достижения точности
нужно пройти порядка $L_1R/\sqrt T$ уровней (с точностью до
логарифмических множителей), а всего уровней $T$. Длина цепочки
выбирается так, чтобы величины $T$ и $L_1R/\sqrt T$ совпадали
с точностью до абсолютных постоянных. Из равенства $T=L_1R/\sqrt T$
получаем $T^{3/2}=L_1R$, то есть $T=(L_1R)^{2/3}$; в
лемме~\ref{lem:lower-Q} выбрано целое $T$ этого порядка.
Точные условия этого выбора учитывают также начальную
невязку и неоднородность. Число обращений к оракулу, необходимое
для поиска следующего звена, оценивается в леммах~\ref{lem:lower-search}
и~\ref{lem:lower-revelations}.

\subsection{Нижняя оценка}

Пусть выполнены предположения~\ref{asm:components}--\ref{asm:accuracy}.
Число $N_0$ задано в~\eqref{eq:complexity-parameters}.

\proofdown{thm:lower}
\begin{theorem}\label{thm:lower}
Для оракула, возвращающего $(f_i(x),\nabla f_i(x))$, выполнено
\[
 \mathcal T^\ast\gtrsim\frac{T_0}{\log^2N_0}.
\]
\end{theorem}

Правая часть не зависит от $\Delta$ при $\Delta\ge2\varepsilon$.
Оценка относится к наихудшей задаче класса в произвольной размерности.
Метод может быть рандомизированным и выбирать точки запросов во всём
пространстве, а не только в линейной оболочке ранее полученных градиентов.

Вместе с верхней оценкой для YETI (теорема~\ref{thm:cheap-upper})
эта теорема показывает, что в данной постановке показатели степеней
при $n$, $L_0$, $L_1$, $R$, $\delta$ и $\varepsilon$ улучшить нельзя
(с точностью до логарифмических множителей). Логарифмический разрыв между верхней и нижней оценками
здесь не устраняется. Предположение~\ref{asm:accuracy} существенно для области
применимости результата: эту нижнюю оценку нельзя автоматически
переносить, например, на подкласс с $\delta=0$. Равенство $\delta=0$ означает,
что точка $x^\ast$ минимизирует каждую компоненту, то есть у всех $f_i$ есть
общая точка минимума. Для этого случая оценки сходимости стохастических
методов $(L_0,L_1)$-SGD и SGD-PS получены в работе
\citet[раздел~7]{gorbunov2024methods}.

\subsection{Доказательство нижней оценки}

\begin{proof}[Схема доказательства теоремы~\ref{thm:lower}]
Задача со скрытой парой функций даёт оценку для $B_n$.
Семейство S даёт оценку для $B_S$, а если его условия
нарушены, используется $B_S\lesssim B_n$.
Семейство X вместе со скрытой парой функций даёт оценку для $B_X$
с точностью до логарифмического множителя. Семейство Q даёт оценку для
$B_Q$. Когда условия~\eqref{eq:lower-Q-regime} не выполнены,
этот член оценивается через уже полученные оценки.

Следовательно,
\[
 T_0\sim\max\{B_n,B_S,B_X,B_Q\}
 \lesssim\mathcal T^\ast\log^2N_0.
\]
Леммы~\ref{lem:lower-finite}--\ref{lem:lower-revelations}
позволяют применить эти оценки к произвольным рандомизированным
методам. Разбор всех режимов параметров приведён
в приложении~\ref{app:lower-constants}.
\end{proof}
